\documentclass[10pt,a4paper]{amsart}
\usepackage[margin=19mm]{geometry}
\usepackage{amsmath,amssymb,mathtools}
\usepackage[scr=rsfs,cal=euler]{mathalfa}
\usepackage{xfrac}
\usepackage[unicode,hidelinks,bookmarks=false]{hyperref}
\newcommand{\ie}{i.e.\ }
\newcommand{\cf}{cf.\ }
\newcommand{\resp}{respectively\ }
\newcommand{\Subsection}{\subsection}
\providecommand{\nobreakdash}{\nobreak\hbox{-}}
\setlength{\emergencystretch}{2em}
\multlinegap0pt
\usepackage{tikz}
\usetikzlibrary{cd}
\usepackage{amssymb}
\usepackage{tikz}
\usepackage[shortlabels]{enumitem}
\usepackage{natbib}
\usepackage{mathtools}
\usepackage[normalem]{ulem}
\newtheorem{lem}{Lemma}
\title{On outer automorphisms of certain graph $\textrm{C}^{\ast}$-algebras}
\author{Swarnendu Datta, Debashish Goswami, Soumalya Joardar}
\date{Source: arXiv:2506.23709v1 (CC BY 4.0)}
\begin{document}
\maketitle

\noindent\emph{Source and licence.} On outer automorphisms of certain graph $\textrm{C}^{\ast}$-algebras. Version arXiv:2506.23709v1 (CC BY 4.0), distributed by arXiv under the Creative Commons Attribution 4.0 International licence, http://creativecommons.org/licenses/by/4.0/, as recorded in the arXiv licence metadata for this exact version and shown on the abstract page https://arxiv.org/abs/2506.23709v1. Full licence notice: \texttt{licenses/arxiv-2506.23709-CC-BY-4.0.txt}.\par\smallskip

\noindent\textit{Editorial scope.} Counted unit: the article's lem lem (source0.tex lines 341--343). The article's complete original proof of it is delivered with it (source0.tex lines 344--365, one \texttt{proof} environment of 22 lines). No mathematics is written, repaired or re-derived: the statement and the proof are the article's own lines, in the article's own notation, and no retained line is substituted or re-flowed.  No further source-local context is needed: every cross-reference of every retained line resolves inside the two delivered spans.  Nothing was omitted from any retained span, so the package carries no omission record.  No retained line contains a citation, so the wrapper carries no bibliography and no bibliography entry.  The retained lines contain the article's own diagram code, which the wrapper typesets with the standard tikz packages; no image file is delivered and the package declares no asset dependency.  Editorial formatting only, in addition: an AMS \texttt{amsart} preamble that restates the article's own declarations and macros used by the retained lines, namely the environments \texttt{lem} on separate counters, together with the standard packages those lines need.  The retained environments print in this document as Lemma 1 in order of appearance, whereas the article prints the same items with the numbering of its own full document.  The complete native source of the article as distributed in the arXiv e-print for this exact version is bundled unchanged as \texttt{sources/180729/source0.tex}.
            \begin{lem}
              An automorphism $\phi$ of a row finite directed graph $G$ descends to an automorphism of the abelian group $K_{0}(C^{\ast}(G))$. We denote the automorphism of $K_{0}(C^{\ast}(G))$ corresponding to a graph automorphism $\phi$ by $\underline{\phi}$.
            \end{lem}

            \begin{proof}
            Recall that $K_{0}(\textrm{C}^{\ast}(G))$ is isomorphic to the cokernel of the group homomorphism $B:\mathbb{Z}V_{E}\rightarrow\mathbb{Z}V$ given on the $\mathbb{Z}$-linear basis $\{v\}_{v\in V_{E}}$ by
            \begin{displaymath}
                B(v)=v-\sum_{e:r(e)=v}s(e).
            \end{displaymath}
            As $\phi\in\textrm {Aut}(G)$, it maps $V_{E}$ to $V_{E}$. Then extending $\phi$ $\mathbb{Z}$-linearly, we get maps from $\mathbb{Z}V_{E}$ to $\mathbb{Z}V_{E}$ and $\mathbb{Z}V$ to $\mathbb{Z}V$. We continue to denote the extensions by $\phi$. For $v\in V_{E}$,
            \begin{displaymath}
                B(\phi(v))=\phi(v)-\sum_{r(e)=\phi(v)}s(e).
            \end{displaymath}
            As $\phi$ is an automorphism of the graph $G$, for any $e\in E$ such that $r(e)=\phi(v)$, $s(e)=\phi(w)$ for some $w\in E$ such that $e^{\prime}=(w,v)\in E$ and  $\phi(e^{\prime})=e$. Consequently, 
            \begin{displaymath}
                B(\phi(v))=\phi\Big(v-\sum_{r(e^{\prime})=v}s(e^{\prime})\Big)=\phi(Bv).
            \end{displaymath}
            Therefore, by the $\mathbb{Z}$-linearity of the maps $\phi$ and $B$, we get the following commutative diagram:
            \[
\begin{tikzcd}
\mathbb{Z}V_{E} \arrow{r}{B} \arrow[swap]{d}{\phi} & \mathbb{Z}V \arrow{d}{\phi} \\
\mathbb{Z}V_{E} \arrow{r}{B} & \mathbb{Z}V
\end{tikzcd}
\]
Hence we get a well-defined group homomorphism $\underline{\phi}:{\rm coker}(B)\rightarrow {\rm coker}(B)$ and consequently a group homomorphism $\underline{\phi}:K_{0}(\textrm{C}^{\ast}(G))\rightarrow K_{0}(\textrm{C}^{\ast}(G))$. Repeating the argument with $\phi^{-1}$, we get a group homomorphism $\underline{\phi}^{-1}:K_{0}(\textrm{C}^{\ast}(G))\rightarrow K_{0}(\textrm{C}^{\ast}(G))$. It is straightforward to verify $(\underline{\phi})^{-1}=\underline{\phi}^{-1}$ proving that $\underline{\phi}\in\textrm {Aut}\Big(K_{0}(\textrm{C}^{\ast}(G))\Big)$.
            \end{proof}

\end{document}
